\par\medskip
\noindent\textbf{Teorema 5.4.}
Sea $X\xrightarrow{f}Y$ un morfismo esencialmente de presentación
finita. Supongamos que existe un entero fijo $n\geq0$ para el cual se
cumple una de las condiciones siguientes:

\smallskip
\noindent\textbf{(1)} $f$ es plano y
$\dim_xX_{f(x)}=n$ para todo $x\in X$;

\smallskip
\noindent\textbf{(2)} $f$ es una intersección completa relativa de
dimensión virtual $n$ (SGA 6, VIII, 1.9).

\smallskip
Entonces $x\notin J_X^n(X|Y)$ si y solo si $f$ es liso en $x$.

\smallskip
\noindent\textit{Demostración.}
Por 5.3(4), $x\notin J_X^n(X|Y)$ si y solo si existe un diagrama
\[
\begin{tikzcd}[column sep=huge,row sep=large]
U \arrow[r,hook,"i"] \arrow[d,"f"']
  & U' \arrow[dl,"f'"] \\
Y & {}
\end{tikzcd}
\]
tal que $f'$ es liso en $x'$, con
$\dim_{x'}U'_{f(x)}=n$, mientras que $i$ es una inmersión cerrada y
$x'=i(x)$. Basta, pues, demostrar que, bajo (1) o (2), $i$ es un
isomorfismo en $x$.

Pongamos
\[
A=\mathcal O_{Y,f(x)},\qquad
B'=\mathcal O_{U',x'},\qquad B=\mathcal O_{U,x},
\]
y sea exacta la sucesión
\[
0\longrightarrow I\longrightarrow B'\longrightarrow B
\longrightarrow0
\]
Debemos demostrar que $I=0$.

Supongamos (1). Como $B$ es plano sobre $A$, la sucesión
\[
0\longrightarrow I/\mathfrak mI
\longrightarrow B'/\mathfrak mB'
\longrightarrow B/\mathfrak mB\longrightarrow0
\]
es exacta, donde $\mathfrak m$ es el ideal maximal de $A$. Por
hipótesis,
\[
\dim_xU_{f(x)}=\dim_{x'}U'_{f(x)},
\]
y, por tanto,
$\dim(B'/\mathfrak mB')=\dim(B/\mathfrak mB)$. Como
$B'/\mathfrak mB'$ es un anillo local regular, se deduce que
$I/\mathfrak mI=0$ y, por tanto, $I=0$.

%%%%%%%%%%  scan idx 129  |  folio 118  %%%%%%%%%%

Supongamos ahora (2), y consideremos la sucesión exacta
\[
I/I^2\xrightarrow{d}B\otimes_{B'}\Omega_{B'|A}
\longrightarrow\Omega_{B|A}\longrightarrow0.
\]
Como $X$ es una intersección completa relativa sobre $Y$ y
$U'\longrightarrow Y$ es liso en $x'$, la aplicación
\[
d:I/I^2\longrightarrow B\otimes_{B'}\Omega_{B'|A}
\]
es inyectiva en $x$, y tanto $I/I^2$ como
$B\otimes_{B'}\Omega_{B'|A}$ son libres en $x$. Por la definición de
la dimensión virtual de una intersección completa relativa, debemos
tener, por tanto,
\[
n=\dim_{\kappa(x)}\Omega_{B'|A}(x')
  -\dim_{\kappa(x)}(I/I^2)(x)
 =n-\dim_{\kappa(x)}(I/I^2)(x).
\]
Así, $I/I^2=0$ y, por tanto, $I=0$. \hfill$\square$
